% SPDX-License-Identifier: Apache-2.0
\section{The Four Real Conflicts}
\label{sec:conflicts}

Everything else the two languages disagree about is spelling, and
Section~\ref{sec:syntax} settles spelling. These four need a decision about
meaning.

\subsection{Does the designer count cycles?}

LHDL says no. TxHDL says yes, and a UART example counts baud intervals
explicitly.

\begin{rulebox}[Both, with a boundary]
The two answers apply to different things, and the conflict disappears once
that is stated. A UART's baud interval is a specification: it comes from the
protocol on the wire, and a compiler that retimed it would break the design.
An adder's latency is not a specification: it comes from the fabric, and a
compiler that fixed it by hand would waste the fabric.
\end{rulebox}

The merged language therefore has two behavioural forms and a rule for
choosing between them.

\begin{itemize}
  \item \term{Timeless form}, spelled \code{flow}, from LHDL. Dataflow and
        tags, no cycle counts. The compiler schedules it, and a cycle count
        written inside one is a compile error.
  \item \term{Sequenced form}, spelled \code{seq}, from TxHDL. Explicit
        ordering with \code{await}. Cycle counts are honoured exactly as
        written.
\end{itemize}

A \code{seq} may call a \code{flow}, and the call absorbs whatever latency
the compiler chose. A \code{flow} may instantiate a \code{seq} behind a
contract, and sees it as a fixed or elastic latency component.

\subsection{Nominal contracts or structural typing?}

LHDL declares conformance with \code{implements}. TxHDL infers it from
matching fields.

\begin{rulebox}[Structural for data, nominal for behaviour]
Each source already uses its mechanism for what that mechanism is good at.
Field compatibility between two transactions is a question about layout, and
structural typing decides layout well. Whether a pipeline satisfies a bus
contract is a question about intent, and intent should be declared, because
two modules can share a field set by accident.
\end{rulebox}

\subsection{Is verification in the language?}

LHDL's twelfth thesis argues for interop with existing testing
infrastructure instead of a non-synthesisable subset. TxHDL puts assertions,
assumptions, cover points and properties in the language.

\begin{rulebox}[Both, split by what the construct talks to]
The thesis argues against a simulation subset: file input and output, string
handling, and stimulus generation. It does not argue against formal
properties. Assertions and properties are read by formal tools and by
synthesis, not by a simulator's runtime library, so they do not create the
subset the thesis objects to.
\end{rulebox}

Stimulus, file handling and scoreboarding move to the foreign function
interface, per LHDL. Assertions, assumptions, cover points and properties
stay in the language, per TxHDL.

\subsection{Where does a pipeline live?}

LHDL makes a pipeline a top-level definition that may implement an
interface, with a \code{pipe} declaration for a value that persists across
stages. TxHDL makes it a member of a module with a call signature and a
return in the last stage.

\begin{rulebox}[TxHDL's shape, LHDL's placement]
A pipeline is an item with a signature. It may be declared at top level and
may implement a contract, which is what LHDL needs for late binding. It may
also be declared inside a module and called, which is what TxHDL needs for
the common case. LHDL's \code{pipe} declaration stays, because live range
analysis needs somewhere to say that a value crosses a stage boundary.
\end{rulebox}
